% =====================================================================
% 8. Resultados
% =====================================================================
\section{Resultados}
\label{sec:resultados}

Salvo que se diga otra cosa, los resultados usan el \textbf{escenario base}:
Google con confianza $\geq 0{,}65$ y Microsoft completo. Las tablas por
municipio están en \ruta{data/processed/resultados_municipio.csv} (680 filas:
170 municipios por cuatro combinaciones de fuente y umbral), y el ranking
completo en \ruta{data/processed/ranking_municipios.csv}.

\subsection{Totales y comparación entre fuentes}

Las dos fuentes casi coinciden en área total (222,17 y
221,13\,km\textsuperscript{2}) aunque Google tiene 1,53 veces más edificios,
porque divide la misma superficie en huellas más pequeñas. La coincidencia
también se da municipio por municipio: la correlación entre el área de las dos
fuentes es de 0,99 (Pearson) y 0,94 (Spearman) sobre los 170 municipios.

\fig[0.85\textwidth]{res_dispersion_municipio.png}{0.4}{Área total por
municipio PDET: Google frente a Microsoft (escala logarítmica). La recta
punteada es 1:1. Los cinco municipios sin edificios de Google no aparecen.}{fig:res-dispersion}

La mayoría de los municipios cae cerca de la recta 1:1. Los que se separan son
pocos, y los más alejados quedan por debajo, es decir, con menos área en
Google: Sipí, López de Micay, Condoto, Carmen del Darién y Ricaurte, todos del
Pacífico.

\subsection{Resultados por subregión}

\tabla{tabla_subregion}

La diferencia entre fuentes no es uniforme. Google tiene más área en el Caribe
(Sierra Nevada - Perijá, Urabá Antioqueño, Montes de María, con $D$ entre
$+0{,}07$ y $+0{,}09$) y Microsoft en el Pacífico y el sur (Chocó, Sur del
Tolima, Pacífico y Frontera Nariñense, con $D$ entre $-0{,}21$ y $-0{,}33$). El
total casi coincide porque los excesos y los faltantes se compensan.

\subsection{Sensibilidad al umbral de confianza}
\label{sec:sensibilidad}

\tabla{sensibilidad}

\fig[0.9\textwidth]{res_sensibilidad_umbral.png}{0.3}{Edificios y área total de
Google según el umbral de confianza, con Microsoft como referencia.}{fig:res-sensibilidad}

Subir el umbral quita muchos más edificios que área: con 0,70 se pierde el
16\,\% de los edificios pero solo el 7\,\% del área, y con 0,75 el 35\,\% y el
19\,\%. Los edificios de baja confianza son, sobre todo, huellas pequeñas. Con
0,75 el área de Google (179,6\,km\textsuperscript{2}) queda 19\,\% por debajo de
Microsoft, así que el umbral sí cambia la conclusión sobre qué fuente tiene más
techo. El orden de los municipios, en cambio, es estable: los 15 con más área
con 0,65 siguen en el top 15 con 0,70, y 14 de 15 con 0,75. Por eso 0,65 es el
escenario base y los otros dos se reportan como sensibilidad.

\subsection{Mapas por municipio}
\label{sec:mapas}

Los mapas (\campo{mapas.py}) usan las mismas clases para las dos fuentes
(quintiles calculados sobre los 340 valores juntos, sin ceros), para que el
mismo color signifique lo mismo en los dos paneles. Las áreas se expresan en
hectáreas (1\,ha $=$ 10\,000\,m\textsuperscript{2}).

\fig{mapa_area_fuentes.png}{0.45}{Área de techos por municipio PDET, Google
(izquierda) y Microsoft (derecha), con las mismas clases. En gris con trama,
los 5 municipios sin edificios en Google.}{fig:mapa-fuentes}

Los dos mapas cuentan la misma historia en casi todo el territorio. El área se
concentra en los municipios con cabeceras urbanas grandes (Santa Marta y
Valledupar superan las 1\,300\,ha con las dos fuentes) y en la Sierra Nevada -
Perijá, Urabá, el norte del Cauca y el piedemonte caqueteño. Los municipios
extensos de la Amazonía y la Orinoquía ocupan buena parte del mapa pero tienen
poca área de techos. La diferencia visible está en el Chocó y la costa
pacífica nariñense, más claros en Google.

\fig[0.9\textwidth]{mapa_diferencia.png}{0.45}{Índice de diferencia $D$ entre
Google y Microsoft por municipio. Naranja: más área en Microsoft; azul: más en
Google; gris: coinciden ($|D| \leq 0{,}1$).}{fig:mapa-diferencia}

En 98 de los 170 municipios las fuentes coinciden dentro de $\pm 10$\,\% y en
otros 49 la diferencia es moderada. Los 23 con desacuerdo fuerte
($|D| > 0{,}3$) no están repartidos al azar: 12 son del Chocó y 6 del Pacífico
y Frontera Nariñense, e incluyen los 5 sin edificios en Google y otros como
Sipí, López de Micay, Ricaurte y Carmen del Darién, donde Google detecta menos
del 10\,\% del área de Microsoft. Solo uno tiene mucha más área en Google:
Istmina ($D = +0{,}60$).

\fig[0.85\textwidth]{mapa_pct_techos.png}{0.42}{Porcentaje de la superficie
municipal cubierta por techos (promedio de las dos fuentes). Clases por
quintiles.}{fig:mapa-pct}

El área total favorece a los municipios grandes; normalizada por la superficie,
la mediana es 0,077\,\%. Piendamó - Tunía (1,33\,\%), Santander de Quilichao
(0,78\,\%), Apartadó (0,65\,\%) y Santa Marta (0,64\,\%) encabezan, y el norte
del Cauca es la zona más densa.

\subsection{Ranking de municipios candidatos}
\label{sec:ranking}

El ranking (\campo{ranking.py}) ordena por área conservadora,
$\min(A_G, A_M)$, y marca los municipios con fuentes en desacuerdo.

\tabla{ranking}

\fig[0.95\textwidth]{ranking_top15.png}{0.4}{Área de techos de los 15
municipios candidatos según cada fuente, ordenados por el área
conservadora.}{fig:ranking}

Los 15 municipios suman 7\,775\,ha de área conservadora, el 39\,\% del total de
los 170 (20\,156\,ha), y ninguno tiene las fuentes en desacuerdo: el criterio
deja fuera, sin necesidad de una regla adicional, a los municipios donde los
datos no son confiables. El ranking es estable: 13 de los 15 coinciden con el
top 15 que saldría de usar solo Google o solo Microsoft, y el umbral de
confianza no cambia cuáles lo encabezan. Están representadas 11 de las 16
subregiones, con Santa Marta, Valledupar y Ciénaga (Sierra Nevada - Perijá) a
la cabeza.

\subsection{Estimación de potencia instalable}
\label{sec:potencial}

El área de techos no es todavía potencial energético. Para dar una idea de la
escala, \campo{potencial.py} convierte el área conservadora de cada municipio
en potencia pico y energía anual con supuestos explícitos, que se leen de
\campo{config.yaml}:

\begin{equation}
  P = A \cdot f_u \cdot \rho, \qquad E = P \cdot H \cdot 365 \cdot PR
  \label{eq:potencial}
\end{equation}

\begin{itemize}
  \item $f_u = 0{,}40$ (rango 0,30 a 0,50): fracción del techo aprovechable
        para módulos, que descuenta sombras, bordes, orientación,
        inclinación y techos que no soportan la carga.
  \item $\rho = 0{,}20$\,kWp/m\textsuperscript{2}: módulos de cerca del
        20\,\% de eficiencia (1\,kW/m\textsuperscript{2} en condiciones
        estándar).
  \item $H = 4{,}5$\,kWh/m\textsuperscript{2}/día: irradiación global
        promedio de Colombia según el Atlas de Radiación Solar de la UPME y
        el IDEAM~\cite{atlas}; se usa un solo valor para todo el país.
  \item $PR = 0{,}80$: rendimiento global del sistema (inversor, temperatura,
        cableado y suciedad).
\end{itemize}

\tabla{potencial}

Con estos supuestos, los 170 municipios PDET tendrían unos 16\,GWp de potencia
instalable en techos (entre 12 y 20\,GWp según el aprovechamiento) y unos
21\,TWh al año, y el top 15 concentra 6,2\,GWp. Es un \textbf{orden de
magnitud}, no un dimensionamiento: es el potencial técnico bruto, que supone
cubrir todos los techos aprovechables, y la irradiación real varía entre
regiones. La sección~\ref{sec:recomendaciones} propone cómo refinarlo.
